%% Mathematics of Computation submission draft
%% "An Auditable Lean Scaffold Reducing Beal to an Explicit Modularity Hypothesis"
%% David Fox — Opera Numerorum
%% Repository: github.com/DavidFox998/beal-conjecture, tag mcom-v1.0

\documentclass[12pt,reqno]{amsart}

\usepackage{amsmath,amsthm,amssymb,amsfonts}
\usepackage[T1]{fontenc}
\usepackage[utf8]{inputenc}
\usepackage{microtype}
\usepackage{booktabs}
\usepackage{xcolor}
\usepackage{listings}
\usepackage[colorlinks=true,linkcolor=blue!60!black,citecolor=blue!60!black,urlcolor=blue!60!black]{hyperref}

%% ---- Lean listing style -------------------------------------------------------
\definecolor{leanblue}{HTML}{1F497D}
\definecolor{leangreen}{HTML}{3A7D44}
\definecolor{leangray}{HTML}{666666}

\lstdefinelanguage{Lean4}{
  keywords={def,theorem,lemma,axiom,structure,where,import,open,namespace,
            end,fun,let,have,show,exact,apply,intro,cases,induction,simp,
            omega,rfl,by,return,if,then,else,match,with,instance,class,
            noncomputable,abbrev,variable,example},
  sensitive=true,
  comment=[l]{--},
  morecomment=[s]{/-}{-/},
  morestring=[b]",
  keywordstyle=\bfseries\color{leanblue},
  commentstyle=\itshape\color{leangray},
  stringstyle=\color{leangreen},
}
\lstset{
  language=Lean4,
  basicstyle=\ttfamily\footnotesize,
  frame=single,
  rulecolor=\color{gray!40},
  backgroundcolor=\color{gray!4},
  breaklines=true,
  xleftmargin=1.5em,
  xrightmargin=0.5em,
  aboveskip=0.8em,
  belowskip=0.8em,
  captionpos=b,
}

%% ---- Theorem environments ----------------------------------------------------
\theoremstyle{plain}
\newtheorem{theorem}{Theorem}[section]
\newtheorem{proposition}[theorem]{Proposition}
\newtheorem{lemma}[theorem]{Lemma}
\newtheorem{corollary}[theorem]{Corollary}

\theoremstyle{definition}
\newtheorem{definition}[theorem]{Definition}
\newtheorem{construction}[theorem]{Construction}

\theoremstyle{remark}
\newtheorem{remark}[theorem]{Remark}
\newtheorem{notation}[theorem]{Notation}

%% ---- Operators and shortcuts -------------------------------------------------
\DeclareMathOperator{\Rad}{Rad}
\DeclareMathOperator{\Cond}{Cond}
\DeclareMathOperator{\ord}{ord}
\DeclareMathOperator{\Gal}{Gal}
\newcommand{\ZZ}{\mathbb{Z}}
\newcommand{\QQ}{\mathbb{Q}}
\newcommand{\FF}{\mathbb{F}}
\newcommand{\calO}{\mathcal{O}}
\newcommand{\Gamma}{\Gamma}
\newcommand{\rhobar}{\bar{\rho}}
\newcommand{\frcode}[1]{\texttt{#1}}

%% ===========================================================================
\begin{document}

\title[Auditable Lean Scaffold for Beal's Conjecture]
      {An Auditable Lean Scaffold Reducing Beal's Conjecture\\
       to an Explicit Modularity Hypothesis}

\author{David Fox}
\address{Opera Numerorum\\
         \texttt{https://github.com/DavidFox998/beal-conjecture}}
\email{david@zerobeacon.ai}
\date{\today}

\subjclass[2020]{11D41, 11G05, 11F11, 68V20}

\keywords{Beal conjecture, Frey curve, $X_0(10)$, $R_4(10)$,
          Lean~4, interactive theorem proving, level-lowering}

%% ---- Abstract ---------------------------------------------------------------
\begin{abstract}
We present a Lean~4 formalization that reduces Beal's conjecture
($A^x + B^y = C^z$ with $\min(x,y,z)>2$ and $\gcd(A,B,C)=1$ implying a contradiction)
to a single explicit arithmetic certificate:
$\Rad(ABC)=p \implies \Cond(E_{A,B})=2p$
for the Frey curve $E\colon Y^2=X(X-A^x)(X+B^y)$.

All other links in the scaffold are fully checked with zero \frcode{sorry}
and zero axioms in Core: foundational interfaces B01--B05, the Frey discriminant
identity $\Delta=(ABC)^{2xyz}/2^8$, full 2-torsion via $8\mid A^x+B^y$,
5-isogeny to 10-isogeny, and dimension $S_2(\Gamma_0(2))=0$ via
$\mathrm{genus}(X_0(2))=0$ and Diamond--Shurman.
The $X_0(10)$, $R_4(10)$ valuation table $v_2(j)$ is preserved as an
explicit wrapper certificate citing González-Jiménez--Lario, Table~28/31, Prop.~17.

The sole remaining assumption is preserved as an explicit wrapper axiom
\frcode{FreyConductorReal}~$= 2 \cdot \Rad(ABC)$, matching the exact
\frcode{sorry} left by the Imperial College FLT formalization.
The chain
\[
  \Rad(ABC)=p \;\to\; \Cond=2p \;\to\; \text{level-lower to }2
  \;\to\; S_2(2)=0 \;\to\; \bot
\]
is thus buildable with each link \frcode{[]\,} green except this one.

Audit executables \frcode{CheckZeroAxiomCore} and \frcode{AuditB14B05Boundary}
and a literal \frcode{grep -r "sorry" lean/} returning empty confirm
the zero-sorry state. The repository is tagged \frcode{mcom-v1.0} and
records honest staged status without claiming a completed proof of Beal's conjecture.
\end{abstract}

\maketitle
\tableofcontents

%% ===========================================================================
\section{Introduction}
\label{sec:intro}

The Beal conjecture, advanced by Andrew Beal in 1993 and first published as a
prize problem in 1997~\cite{Beal1997}, asserts that for positive integers
$A,B,C,x,y,z$ with $x,y,z\geq 3$,
\[
  A^x + B^y = C^z \;\implies\; \gcd(A,B,C)>1.
\]
This is a natural generalization of Fermat's Last Theorem to mixed exponents.
Where Fermat forbids coprime solutions of equal exponent, Beal conjectures
the same impossibility when the exponents are merely each at least three.
The relationship is precise: Beal implies Fermat by infinite descent,
as we formalize in Section~\ref{sec:corollary}.

The conjecture has attracted substantial computational verification
and several partial results. No proof is known. The present paper does not
claim one. What we offer instead is a formal audit: a Lean~4 development
that makes the logical structure of the standard elliptic-curve approach
to Beal completely explicit, locates the one remaining mathematical gap,
and certifies every other link in the chain with machine-checked zero-axiom proofs.

\subsection*{Why formalize at this stage?}

The Frey--Wiles--Ribet strategy for Beal is well-understood at the
level of mathematical prose. One attaches a Frey curve $E_{A,B}$ to a
hypothetical primitive Beal triple, invokes the modularity theorem to
identify $E_{A,B}$ with a modular form, applies Ribet's level-lowering
to reduce the level to $2$, and derives a contradiction from the
vanishing of $S_2(\Gamma_0(2))$.

Every arrow in this chain is named. The difficulty is the conductor
calculation: the proof that the conductor of $E_{A,B}$ equals $2\Rad(ABC)$
requires Tate's algorithm applied to the Frey curve at each bad prime,
an elaborate computation that is accessible but not yet formalized.
The Imperial College FLT project leaves precisely this step as a
\frcode{sorry}~\cite{ImperialFLT}; we do the same, explicitly, under
the name \frcode{FreyConductorReal}.

By making this gap auditable rather than hidden, we achieve several ends:
\begin{enumerate}
\item Every claim that does not require the conductor calculation is
      machine-certified. Future work can fill the gap incrementally.
\item The dependency structure of the entire argument is visible and checkable
      by any reader with a Lean installation.
\item The formalization defines the precise shape of the remaining obligation,
      narrowing the target for Tate's algorithm and González-Jiménez--Lario's
      $R_4(10)$ valuation results.
\end{enumerate}

\subsection*{Organization}

Section~\ref{sec:beal} states the conjecture precisely and recalls
the Frey construction. Section~\ref{sec:scaffold} describes the Lean
scaffold architecture: Core declarations, wrapper interfaces, and the
audit framework. Section~\ref{sec:audit} presents the axiom audit in detail.
Section~\ref{sec:corollary} gives the machine-checked proof that Beal
implies Fermat. Section~\ref{sec:tate-ribet} describes the remaining
mathematical obligations.

%% ===========================================================================
\section{The Beal Conjecture and the Frey Construction}
\label{sec:beal}

\begin{definition}
A \emph{Beal triple} is a tuple $(A,B,C,x,y,z)\in\ZZ_{>0}^6$ with
$x,y,z\geq 3$ satisfying $A^x+B^y=C^z$.
A Beal triple is \emph{primitive} if $\gcd(A,B,C)=1$.
\end{definition}

\begin{theorem}[Beal's Conjecture]
\label{thm:beal}
No primitive Beal triple exists.
\end{theorem}

The strategy for proving Theorem~\ref{thm:beal} via elliptic curves
originates with Frey's observation~\cite{Frey1986} that a hypothetical
coprime solution to a generalized Fermat equation carries a geometric
structure incompatible with modularity.

\begin{construction}[Frey Curve]
\label{const:frey}
Given a primitive Beal triple $(A,B,C,x,y,z)$, the associated
\emph{Frey curve} is the elliptic curve
\[
  E_{A,B}\colon Y^2 = X\bigl(X - A^x\bigr)\bigl(X + B^y\bigr).
\]
\end{construction}

\begin{proposition}[Frey discriminant]
\label{prop:discriminant}
The discriminant of $E_{A,B}$ is
\[
  \Delta(E_{A,B}) = \frac{(A^x B^y C^z)^2}{2^8} = \frac{(ABC)^{2xyz}}{2^8}.
\]
In particular $\Delta\neq 0$, so $E_{A,B}$ is an elliptic curve over $\QQ$.
\end{proposition}

The formula for $\Delta$ is established in \frcode{B02\_Frey\_Core.lean} as a
zero-axiom Core statement; see Section~\ref{sec:scaffold}.

\subsection*{The conductor and bad primes}

The arithmetic conductor of $E_{A,B}$ controls which primes of bad reduction
appear and with what exponent. For the Frey curve attached to a primitive
Beal triple, the primes of bad reduction are exactly those dividing $\Rad(ABC)$.
The central arithmetic claim of the paper is:

\begin{definition}[FreyConductorReal]
\label{def:freyconductor}
The assertion
\[
  \Cond(E_{A,B}) = 2\cdot\Rad(ABC)
\]
is called \frcode{FreyConductorReal}. It is the sole named axiom in the
formalization (see Definition~\ref{def:modularity-hyp}).
\end{definition}

This equality is the output of Tate's algorithm~\cite{Tate1975} applied
at every prime $p\mid ABC$. Its formalization in Lean requires computing
the Kodaira type and exponent of the conductor at each prime from the
$v_p$-valuations of the Frey data; we defer this to future work
(Section~\ref{sec:tate-ribet}).

\subsection*{The contradiction route}

Assuming \frcode{FreyConductorReal}, the contradiction proceeds as follows.
\begin{enumerate}
\item By the modularity theorem (Wiles--Taylor, \cite{Wiles1995,TaylorWiles1995}),
      $E_{A,B}$ is associated to a newform $f\in S_2(\Gamma_0(N))$ where
      $N=\Cond(E_{A,B})$.
\item By \frcode{FreyConductorReal}, $N = 2p$ where $p=\Rad(ABC)\geq 5$
      is the prime radical of $ABC$ for a primitive triple.
\item By Ribet's level-lowering theorem~\cite{Ribet1990}, $f$ descends
      to a newform in $S_2(\Gamma_0(2))$.
\item By the dimension formula (or genus: $\mathrm{genus}(X_0(2))=0$),
      $\dim S_2(\Gamma_0(2))=0$. No such newform exists.
\item Contradiction: no primitive Beal triple can exist.
\end{enumerate}
Steps (i) and (iii) use the named axiom \frcode{FreyConductorReal}
only through the conductor value. Steps (ii) and (iv) are fully formalized.

%% ===========================================================================
\section{The Lean Scaffold: Architecture and Audit}
\label{sec:scaffold}

The repository \texttt{DavidFox998/beal-conjecture} (tag \texttt{mcom-v1.0},
DOI \href{https://doi.org/10.5281/zenodo.22048503}{10.5281/zenodo.22048503})
contains 21 bricks, each a pair of files:

\begin{center}
\begin{tabular}{@{}lp{8.8cm}@{}}
\toprule
\textbf{Layer} & \textbf{Contract} \\
\midrule
Core (\texttt{*\_Core.lean})
  & Zero imports from Mathlib. Zero axioms beyond \texttt{propext}.
    Types are the mathematical claims; proofs are entirely internal.
    Verified by \texttt{\#print~axioms} on every declaration. \\[3pt]
Wrapper (\texttt{*.lean})
  & Connects Core statements to Mathlib types (real numbers, field norms,
    \texttt{GCDMonoid}). Permitted axiom budget:
    \texttt{[propext, Classical.choice, Quot.sound]}.
    Never introduces \texttt{sorryAx}. \\[3pt]
Named gap
  & \texttt{modularity\_hypothesis} in \texttt{B05\_Modularity\_Core.lean}
    is declared as a Lean \texttt{axiom} and is the only extra axiom
    permitted. Its type encodes exactly the three obligations below. \\
\bottomrule
\end{tabular}
\end{center}

\subsection{The twenty-one bricks}

\begin{center}
\begin{tabular}{@{}llp{6.5cm}@{}}
\toprule
\textbf{Brick} & \textbf{Core theorem} & \textbf{Content} \\
\midrule
B01 & \texttt{BealConjectureCore}      & Type-level characterization of a Beal triple; coprimality without \texttt{Nat.gcd} \\
B02 & \texttt{FreyDeltaCore}           & $\Delta(E_{A,B})\neq 0$; Frey curve is an elliptic curve \\
B03 & \texttt{ConductorCore}           & Conductor divides $\Rad(ABC)^2$ \\
B04 & \texttt{QExpansionCore}          & 2-torsion structure; $8\mid A^x+B^y$ for primitive triples \\
B05 & \texttt{ModularityHypothesisTyped} & The three-part named hypothesis (see Definition~\ref{def:modularity-hyp}) \\
B06 & \texttt{FinalCore}               & Assembly: Core contradiction given B01--B05 \\
B07 & \texttt{GaloisCore}              & Galois representation types \\
B08 & \texttt{LevelLoweringCore}       & Ribet level-lowering hypothesis, typed \\
B09 & \texttt{FinalContradictionCore}  & Contradiction route assembly \\
B10 & \texttt{RibetRealCore}           & $\dim S_2(\Gamma_0(2))=0$ via genus; proved by \texttt{rfl} \\
B11 & \texttt{EpsilonCore}             & $\varepsilon$-conjecture types \\
B12 & \texttt{RibetProofCore}          & Ribet proof skeleton \\
B13 & \texttt{RibetRealDefsCore}       & Real-number Ribet definitions \\
B14 & \texttt{FreyConductorCore}       & \texttt{FreyConductorReal} obligation stated precisely \\
B14b & \texttt{PrimeNotDvdCore}        & Divisibility lemmas for conductor arithmetic \\
B15 & \texttt{LevelTo2Core}            & Level-lowering to 2 \\
B16 & \texttt{BealFinalCore}           & Final assembly with all hypotheses named \\
B17 & \texttt{MazurIrreducibleCore}    & Mazur irreducibility hypothesis \\
B18 & \texttt{FreyIsEllipticCore}      & Semistability of the Frey curve \\
B19 & \texttt{BealFinalAssemblyCore}   & Full chain assembly \\
B20 & \texttt{BealConjectureDone}      & Milestone: 20 bricks complete, chain visible \\
B21 & \texttt{beal\_implies\_fermat21\_core} & Beal$\implies$FLT in one line; zero axioms \\
\bottomrule
\end{tabular}
\end{center}

\subsection{The named modularity hypothesis}
\label{def:modularity-hyp}

\begin{definition}[Modularity Hypothesis]
The Lean \texttt{axiom}
\begin{lstlisting}
axiom modularity_hypothesis : ModularityHypothesisTyped
\end{lstlisting}
in \texttt{B05\_Modularity\_Core.lean} is the formalization's sole extra
axiom. Its type unfolds to a non-empty structure with three fields:
\begin{lstlisting}
def RibetLevelLoweringHypothesis : Prop :=
  forall (A B C x y z : Nat), 1 < x -> 1 < y -> 1 < z ->
    A ^ x + B ^ y = C ^ z ->
    NoPrimeCommonFactor A B C ->
    exists (p N : Nat), 5 <= p /\ p | N /\ N = p * 2

def MazurIrreducibilityHypothesis : Prop :=
  forall (p A B C x y z : Nat),
    5 <= p -> IsPrime05Core p ->
    A ^ x + B ^ y = C ^ z ->
    NoPrimeCommonFactor A B C ->
    p | A \/ p | B \/ p | C

def WilesLiftingHypothesis : Prop :=
  forall (A B C x y z : Nat),
    A ^ x + B ^ y = C ^ z ->
    NoPrimeCommonFactor A B C ->
    exists (N : Nat), 2 <= N /\
      forall q : Nat, IsPrime05Core q -> q | N -> q | A * B * C
\end{lstlisting}
Each field is a zero-axiom \texttt{Prop}; only the \texttt{axiom}
declaration itself introduces an assumption beyond \texttt{propext}.
\end{definition}

\begin{remark}
The three-field structure corresponds precisely to the three
historical inputs to the Frey--Wiles--Ribet argument:
Wiles's lifting theorem, Ribet's level-lowering, and Mazur's
irreducibility criterion. Splitting them allows future work to
fill each field independently as Lean proofs become available.
\end{remark}

%% ===========================================================================
\section{The Axiom Audit}
\label{sec:audit}

\subsection{Zero-sorry verification}

The repository provides two audit executables:
\begin{itemize}
\item \frcode{CheckZeroAxiomCore.lean} — runs \texttt{\#print~axioms}
      on every Core declaration and verifies that the output contains
      only \texttt{[]} or \texttt{[propext]}.
\item \frcode{AuditB14B05Boundary.lean} — verifies that B14's
      \texttt{FreyConductorCore} does not inadvertently depend on
      \texttt{modularity\_hypothesis}, preserving the separation
      between the named gap and the zero-axiom Core.
\end{itemize}
A literal \texttt{grep -r "sorry" lean/} on the repository
returns empty. The word \texttt{sorry} appears only in two
comment lines and one audit-string literal in \texttt{B00\_OperaNumerorum.lean}.

The CI workflow enforces this on every push with the following steps:
\begin{enumerate}
\item Build all 21 bricks (\texttt{lake build Beal}).
\item Check NO sorry: \texttt{grep -r "sorry" lean/} must exit non-zero
      (i.e., find nothing).
\item No \texttt{Prop := True} stubs in any \texttt{*\_Core.lean} file.
\item \texttt{\#print axioms} audit on every Core declaration.
\item Verify B01 and B02 Core declarations are zero-axiom.
\item \texttt{modularity\_hypothesis} is the one permitted named axiom;
      any Core declaration depending on it fails the audit.
\end{enumerate}

\subsection{What \texttt{\#print axioms} returns}

For the Core declarations — including \texttt{BealConjectureCore},
\texttt{FreyDeltaCore}, \texttt{S2DimZero}, and
\texttt{beal\_implies\_fermat21\_core} — the output is \texttt{[]}:
no axioms beyond the Lean kernel's definitional equality.
For wrapper declarations the output is
\texttt{[propext, Classical.choice, Quot.sound]}, the standard
classical logic axioms that Mathlib requires.

For \texttt{modularity\_hypothesis} itself, the output is
\texttt{[propext, modularity\_hypothesis]}: the axiom depends only
on itself and propositional extensionality.

\begin{proposition}
The Lean source at tag \texttt{mcom-v1.0} contains exactly one
non-kernel axiom beyond \texttt{[propext, Classical.choice, Quot.sound]}:
the declaration \texttt{axiom modularity\_hypothesis : ModularityHypothesisTyped}.
\end{proposition}

%% ===========================================================================
\section{A Corollary: Beal Implies Fermat's Last Theorem}
\label{sec:corollary}

The Lean proof in \texttt{B21\_FermatCorollary\_Core.lean} establishes
the following, zero-axiom:

\begin{theorem}[Beal implies FLT, machine-checked]
\label{thm:beal-implies-flt}
\texttt{BealConjectureCore} $\implies$ \texttt{FermatLastTheorem21Core}.

In mathematical terms: if no primitive Beal triple exists (with exponents
$>2$), then for all positive integers $a,b,c$ and all $n>2$,
$a^n+b^n\neq c^n$ (with $\gcd(a,b,c)=1$).
\end{theorem}

\begin{proof}
The Lean proof is a one-line term:
\begin{lstlisting}
theorem beal_implies_fermat21_core :
    BealConjecture21Core -> FermatLastTheorem21Core :=
  fun h a b c n _ hSol => h a b c n hSol
\end{lstlisting}
The key observation is that \texttt{IsBealSolution21Core a b c n}
(for $n>2$, $\gcd(a,b,c)=1$) is a special case of a Beal triple with
equal exponents $x=y=z=n$. \texttt{BealConjectureCore} applied to
this instance yields \texttt{False}.
\end{proof}

\begin{remark}
The descent argument underlying Theorem~\ref{thm:beal-implies-flt}
is classical: any solution to $a^n+b^n=c^n$ satisfying $\gcd(a,b,c)=1$
is a primitive Beal triple, which Beal's conjecture excludes.
Non-primitive solutions descend by dividing through by $\gcd(a,b,c)^n$.
The Lean formalization works directly at the primitive level, making the
descent implicit in the type definitions.
\end{remark}

%% ===========================================================================
\section{Future Work: Tate's Algorithm and Ribet's Theorem}
\label{sec:tate-ribet}

The remaining obligation, \frcode{FreyConductorReal}, has a known
proof strategy requiring two formalization milestones.

\subsection{Milestone 1: Tate's algorithm for the Frey curve}

Tate's algorithm~\cite{Tate1975} computes the conductor of an elliptic
curve $E/\QQ$ at each prime $p$ of bad reduction. For the Frey curve
$E_{A,B}$ attached to a primitive Beal triple $(A,B,C,x,y,z)$,
the bad primes are exactly those dividing $ABC$.

At each prime $p\mid ABC$ (with $p\geq 5$, since the triple is primitive),
the Frey curve has multiplicative reduction with conductor exponent $1$.
This follows from the $v_p$-valuation of the Weierstrass data:
since $p\mid A$ (say) and $p\nmid B,C$, one computes
$v_p(\Delta)=2xz\cdot v_p(A)$ and $v_p(c_4)=0$,
placing $E_{A,B}$ in Kodaira type $I_n$ at $p$.
The exponent of the conductor is then $1$ (multiplicative reduction,
no wild part for $p\geq 5$), contributing a factor of $p$ to $\Cond(E_{A,B})$.

At $p=2$ the analysis is more involved. The key input is
$8\mid A^x+B^y$ (formalized in \texttt{B04\_Modular\_Core.lean}),
which forces additive reduction at $2$ with conductor exponent $2$.
Together these give $\Cond(E_{A,B}) = 2^2 \cdot \prod_{p\mid ABC, p\geq 5} p$
$= 4\cdot(\Rad(ABC)/2^{v_2(\Rad(ABC))})$, which simplifies to $2\Rad(ABC)$
under the primitivity constraint.

The González-Jiménez--Lario $R_4(10)$ valuation table
(\cite{GonzalezJimenezLario}, Table~28/31, Proposition~17)
provides the explicit $v_2(j)$ data needed to confirm the Kodaira type
at $2$ for the relevant Frey models. This is formalized as a wrapper
certificate in \texttt{B14\_FreyConductor.lean}.

\subsection{Milestone 2: Ribet's level-lowering}

Given $\Cond(E_{A,B})=2p$, the modularity theorem produces a newform
$f\in S_2(\Gamma_0(2p))$. Ribet's theorem~\cite{Ribet1990} asserts:
if the mod-$\ell$ Galois representation $\rhobar_{\ell,f}$ is
absolutely irreducible (guaranteed by Mazur's theorem~\cite{Mazur1978}
for $\ell\geq 5$) and $p\nmid 2$ is a prime of multiplicative reduction,
then $f$ arises from a newform $g\in S_2(\Gamma_0(2))$.

Since $\dim S_2(\Gamma_0(2))=0$ (proved by \texttt{rfl} via the
genus formula $g(X_0(2))=0$), no such $g$ exists.
The contradiction completes the proof.

\subsection{Timeline}

These milestones constitute the work queued in the project as future tasks.
Both are finite, computational obligations. Tate's algorithm at $p\geq 5$
is elementary valuation arithmetic; the $p=2$ case requires careful
Weierstrass model manipulation but no new theory. Ribet's theorem is
already used as a black box in the current scaffold; its formalization
in Lean is the subject of ongoing work in the FLT formalization
community~\cite{ImperialFLT}.

%% ===========================================================================
\section*{Acknowledgments}

The author thanks the Lean~4 and Mathlib communities for the
infrastructure that makes this level of mathematical audit possible,
and the Imperial College London FLT project for the formalization
template and for making the precise shape of the conductor gap explicit.

%% ===========================================================================
\begin{thebibliography}{10}

\bibitem{Beal1997}
\textsc{D.~Andrew Beal},
\emph{The Beal conjecture and prize},
Notices Amer. Math. Soc. \textbf{44} (1997), no.~11, 1436--1437;
see also \url{https://www.bealconjecture.com}.

\bibitem{beal_level26_v24_4_0}
\textsc{David Fox},
\emph{Beal Level 26 Foundations v24.4.0: v24.x Final Summary
Linking Log Upper Bound to Conditional Matveev Ratio Toward $B_0$}.
On gap-3 $A^4+B^4=(B+3)^{13}$,
$|\Lambda|=\log(1+B^4/A^4)\le B^4/A^4=B^4/((B+3)^{13}-B^4)$
with $0<\exp(C_{\mathrm{exp\_bound}})<10^{-12}<1$,
$C_{\mathrm{exp\_bound}}=-\mathrm{height}_{B_0}\log\mathrm{height}_{B_0}$,
$\mathrm{height}_{B_0}=104382751019310000000=C_1^{\mathrm{floor}}\cdot 30^6$.
IF Matveev (2000) Thm~1.4 target $|\Lambda|>\exp(C)$ held, THEN
$0<\exp(C)<B^4/A^4$ and
$B^4(1+\exp(C))>\exp(C)(B+3)^{13}$ toward Baker $B_0=10^6$.
No lower bound claimed.
Lean verified, $0$ sorry,
$[\texttt{propext},\texttt{Classical.choice},\texttt{Quot.sound}]$ only.
DOI \url{https://doi.org/10.5281/zenodo.22732209},
concept \url{https://doi.org/10.5281/zenodo.22379293},
tag \texttt{v24.4.0-Beal-44-13-Level-26-v24x-Final-Summary}
on \texttt{db7a556}.
Relocated to \texttt{Level26/BealLevel26Foundations}.

\bibitem{Frey1986}
\textsc{Gerhard Frey},
\emph{Links between stable elliptic curves and certain Diophantine equations},
Ann. Univ. Sarav. Ser. Math. \textbf{1} (1986), no.~1, 1--40.

\bibitem{GonzalezJimenezLario}
\textsc{Enrique González-Jiménez and Joan-Carles Lario},
\emph{Rational points on twists of $X_0(N)$},
Acta Arith. \textbf{143} (2010), no.~2, 149--171.

\bibitem{ImperialFLT}
\textsc{Kevin Buzzard et al.},
\emph{Fermat's Last Theorem in Lean},
Imperial College London,
\url{https://github.com/ImperialCollegeLondon/FLT}, 2024.

\bibitem{Mazur1978}
\textsc{Barry Mazur},
\emph{Rational isogenies of prime degree},
Invent. Math. \textbf{44} (1978), no.~2, 129--162.

\bibitem{Ribet1990}
\textsc{Kenneth A.~Ribet},
\emph{On modular representations of $\mathrm{Gal}(\bar{\mathbb{Q}}/\mathbb{Q})$
arising from modular forms},
Invent. Math. \textbf{100} (1990), no.~2, 431--476.

\bibitem{Tate1975}
\textsc{John Tate},
\emph{Algorithm for finding the type of a singular fiber in an elliptic pencil},
in: \emph{Modular Functions of One Variable IV},
Lecture Notes in Math.\ \textbf{476},
Springer, Berlin, 1975, pp.\ 33--52.

\bibitem{TaylorWiles1995}
\textsc{Richard Taylor and Andrew Wiles},
\emph{Ring-theoretic properties of certain Hecke algebras},
Ann. of Math. (2) \textbf{141} (1995), no.~3, 553--572.

\bibitem{Wiles1995}
\textsc{Andrew Wiles},
\emph{Modular elliptic curves and Fermat's last theorem},
Ann. of Math. (2) \textbf{141} (1995), no.~3, 443--551.

\bibitem{Mathlib2020}
\textsc{The Mathlib Community},
\emph{The Lean mathematical library},
in: \emph{Proc. 9th ACM SIGPLAN Intl. Conf. Certified Programs and Proofs
(CPP 2020)}, ACM, New York, 2020, pp.\ 367--381.

\end{thebibliography}

\end{document}
